% =============================================================================
% Section 3.4: Intent-to-Workflow Decision Accuracy
% Formatted for two-column Elsevier layout
% Wide tables use table* to span both columns
% =============================================================================

\subsection{Intent-to-Workflow Decision Accuracy}

To evaluate whether \dentalclaw{} can reliably translate natural-language research
intents into appropriate workflows, we constructed a 6-dimension weighted evaluation
framework informed by established agent evaluation benchmarks
(OpenAI~Eval~\cite{openai2024}, TRAJECT-Bench~\cite{traject2025},
TRACE~\cite{trace2025}, Claw-Eval, OpenClawBench, and WildClawBench) and tested it
on 11 representative intents spanning four categories: standard executable workflows,
boundary cases, trap queries, and ambiguous requests.

The six dimensions and their associated weights are listed in Table~\ref{tab:eval_dims}.
Each intent was processed through the full platform pipeline
(language~detection$\rightarrow$parsing$\rightarrow$academic~literature~search$\rightarrow$agent~decision$\rightarrow$method~selection).
The agent decision layer used OpenClaw with a DeepSeek~V4~Pro fallback; academic
search drew from Europe~PMC and arXiv. Each run took approximately 50--80~s.

\begin{table*}[!t]
\centering
\caption{Six-dimension weighted evaluation framework. Weights reflect the relative
importance of each dimension. Reference sources indicate the established benchmarks
from which each dimension was adapted.}
\label{tab:eval_dims}
\footnotesize
\begin{tabular}{@{}p{3.0cm} c p{5.5cm} p{4.5cm}@{}}
\hline
\textbf{Dimension} & \textbf{Weight} & \textbf{Evaluated Property} & \textbf{Reference Source(s)} \\
\hline
Planning   & 0.35 & \texttt{supported}/\texttt{executable}/\texttt{method} correctness & Claw-Eval, OpenAI Eval \\
QC Blocking & 0.20 & Correct reject vs.\ allow decisions  & TRACE, WildClawBench \\
Intent Parsing & 0.15 & \texttt{dataset}/\texttt{task}/\texttt{mode} extraction  & OpenAI Eval \\
External Proposal & 0.15 & Agent proposal quality when no registry match & OpenClawBench \\
Ambiguity Handling & 0.10 & Correct rejection of underspecified intents  & WildClawBench \\
Boundary Identification & 0.05 & Correct rejection of out-of-scope requests  & TRAJECT-Bench, WildClawBench \\
\hline
\end{tabular}
\end{table*}

Table~\ref{tab:intent_results} reports the decision accuracy across the 11 intents.
The platform correctly selected the registered method for all five standard intents
(1.00), correctly rejected both trap intents (1.00), and correctly asked for
clarification on ambiguous and boundary intents where the task or dataset was not
fully specified (1.00). For the two boundary intents involving training and detection,
no registered method existed; the agent proposed external solutions from the
codebase---an nnU-Net auto-training skill and a YOLO detection script---with confidence
scores of 0.85 and 0.70 respectively, flagged as \texttt{external\_suggestion}
requiring human review (0.98).

\begin{table*}[!t]
\centering
\caption{Intent-to-workflow decision accuracy across 11 representative intents.
Weighted average across all intents using the dimension weights from
Table~\ref{tab:eval_dims}.}
\label{tab:intent_results}
\footnotesize
\begin{tabular}{@{}c p{5.5cm} p{2.0cm} c p{4.0cm}@{}}
\hline
\# & \textbf{Intent} & \textbf{Category} & \textbf{Score} & \textbf{Method Selected} \\
\hline
 1 & TDD panoramic tooth segmentation inference   & standard  & 1.00 & \texttt{tdd\_2d\_seg\_infer\_report} \\
 2 & TDD panoramic super-resolution enhancement    & standard  & 1.00 & \texttt{dental\_2d\_super\_resolution} \\
 3 & ToothFairy3 CBCT 3D tooth segmentation        & standard  & 1.00 & \texttt{toothfairy3\_3d\_seg\_...} \\
 4 & Private dental data segmentation training     & standard  & 1.00 & \texttt{private\_2d\_segmentation\_train} \\
 5 & Dental anomaly detection (panoramic)          & standard  & 1.00 & \texttt{dental\_anomaly\_detection} \\
 6 & ``Build a model with TDD'' (vague task)       & ambiguous & 1.00 & \texttt{platform\_clarification} \\
 7 & ``Write a dental AI paper'' (non-CV)          & trap      & 1.00 & Rejected (\texttt{supported=false}) \\
 8 & ``Video action recognition'' (non-dental)     & trap      & 1.00 & Rejected (\texttt{supported=false}) \\
 9 & Inference on private data (no task specified) & boundary  & 1.00 & \texttt{platform\_clarification} \\
10 & TDD tooth segmentation \emph{training}        & boundary  & 0.98 & \texttt{agent\_external\_suggestion} \\
11 & TDD bbox tooth \emph{detection} training      & boundary  & 0.98 & \texttt{agent\_external\_suggestion} \\
\hline
\multicolumn{5}{@{}c@{}}{Weighted average: \textbf{0.977} \quad (all $\ge$~0.70 threshold)} \\
\end{tabular}
\end{table*}

Table~\ref{tab:dim_scores} shows the per-dimension breakdown. All five standard
intents matched their intended registry methods exactly, both trap intents were
correctly rejected, and both boundary intents were correctly flagged as requiring
human review. No incorrect workflow was ever silently executed.

\begin{table}[!t]
\centering
\caption{Per-dimension scores (averaged over 11 intents).}
\label{tab:dim_scores}
\footnotesize
\begin{tabular}{@{}l c c@{}}
\hline
\textbf{Dimension} & \textbf{Weight} & \textbf{Score} \\
\hline
Planning               & 0.35 & 0.955 \\
QC Blocking            & 0.20 & 1.000 \\
Intent Parsing         & 0.15 & 0.970 \\
External Proposal      & 0.15 & 1.000 \\
Ambiguity Handling     & 0.10 & 1.000 \\
Boundary Identification& 0.05 & 0.945 \\
\hline
\textbf{Weighted Total}& 1.00 & \textbf{0.977} \\
\hline
\end{tabular}
\end{table}

These results provide quantitative evidence that \dentalclaw's{} agent-driven
pipeline can correctly interpret diverse research intents and select appropriate
workflows, with reliable rejection of out-of-scope requests and appropriate
flagging of novel but in-scope proposals for human review.
